
\documentclass{article}
\usepackage{colortbl}
\usepackage{makecell}
\usepackage{multirow}
\usepackage{supertabular}

\begin{document}

\newcounter{utterance}

\centering \large Interaction Transcript for game `wordle', experiment `high\_frequency\_words\_no\_clue\_no\_critic', episode 10 with nemotron{-}3{-}ultra.
\vspace{24pt}

{ \footnotesize  \setcounter{utterance}{1}
\setlength{\tabcolsep}{0pt}
\begin{supertabular}{c@{$\;$}|p{.15\linewidth}@{}p{.15\linewidth}p{.15\linewidth}p{.15\linewidth}p{.15\linewidth}p{.15\linewidth}}
    \# & \multicolumn{2}{c}{Player} && \multicolumn{2}{c}{Game Master} \\
    \hline

    \theutterance \stepcounter{utterance}  
    & & & \multicolumn{4}{p{0.6\linewidth}}{
        \cellcolor[rgb]{0.9,0.9,0.9}{
            \makecell[{{p{\linewidth}}}]{
                \texttt{\tiny{[P1$\langle$GM]}}
                \texttt{Eres un mago del lenguaje al que le gusta adivinar palabras usando las reglas dadas.} \\
\\ 
\texttt{¡Bienvenido a Wordle! Tienes seis intentos para adivinar la palabra objetivo, una palabra válida en español de cinco letras minúsculas. Por favor, usa las etiquetas "explicación:" e "intento:" para proporcionar una explicación concisa para cada intento.} \\
\\ 
\texttt{Por ejemplo, si tu intento es "parte", tu respuesta debería ser} \\
\texttt{explicación: esta es una palabra común en español de cinco letras, y estoy comenzando mi intento con esta palabra.} \\
\texttt{intento: parte} \\
\\ 
\texttt{Después de cada intento, tu respuesta será validada y recibirás retroalimentación indicando qué letras son correctas (verde), qué letras son correctas pero están en la posición incorrecta (amarillo) y qué letras son incorrectas (rojo). Esta retroalimentación puede ser útil para determinar qué letras incluir o excluir en tu siguiente intento.} \\
\\ 
\texttt{Por ejemplo, la retroalimentación para "parte" podría ser:} \\
\texttt{retroalimentación\_intento: p<amarillo> a<amarillo> r<verde> t<amarillo> e<rojo>} \\
\\ 
\texttt{La explicación debe contener detalles sobre cómo se usa el retroalimentación\_intento para llegar a un nuevo intento.} \\
\\ 
\texttt{No generes ningún otro texto.} \\
\\ 
\texttt{Comencemos con tu primer intento.} \\
            }
        }
    }
     \\ \\

    \theutterance \stepcounter{utterance}  
    & \multicolumn{4}{p{0.6\linewidth}}{
        \cellcolor[rgb]{0.561,0.737,0.561}{
            \makecell[{{p{\linewidth}}}]{
                \texttt{\tiny{[P1$\rangle$GM]}}
                \texttt{explicación: comienzo con una palabra que contiene vocales frecuentes (a, e, o) y consonantes muy comunes en español (r, s) para maximizar la información inicial.} \\
\texttt{intento: areso} \\
            }
        }
    }
    & & \\ \\

    \theutterance \stepcounter{utterance}  
    & & & \multicolumn{2}{p{0.3\linewidth}}{
        \cellcolor[rgb]{0.9,0.9,0.9}{
            \makecell[{{p{\linewidth}}}]{
                \texttt{\tiny{[GM$|$GM]}}
                \texttt{Error: The guessed word is not a valid word for this game.} \\
            }
        }
    }
    & & \\ \\

    \theutterance \stepcounter{utterance}  
    & & & \multicolumn{4}{p{0.6\linewidth}}{
        \cellcolor[rgb]{0.9,0.9,0.9}{
            \makecell[{{p{\linewidth}}}]{
                \texttt{\tiny{[P1$\langle$GM]}}
                \texttt{Tu intento no es una palabra válida para este juego. Por favor, inténtalo de nuevo.} \\
\\ 
\texttt{Proporciona tu respuesta solo en este formato.} \\
\texttt{explicación: detalles} \\
\texttt{intento: palabra} \\
            }
        }
    }
     \\ \\

    \theutterance \stepcounter{utterance}  
    & \multicolumn{4}{p{0.6\linewidth}}{
        \cellcolor[rgb]{0.561,0.737,0.561}{
            \makecell[{{p{\linewidth}}}]{
                \texttt{\tiny{[P1$\rangle$GM]}}
                \texttt{explicación: "areso" no es válida; pruebo "arroz", palabra común con vocales abiertas y consonantes frecuentes (r, z) para cubrir patrones típicos.} \\
\texttt{intento: arroz} \\
            }
        }
    }
    & & \\ \\

    \theutterance \stepcounter{utterance}  
    & & & \multicolumn{2}{p{0.3\linewidth}}{
        \cellcolor[rgb]{0.9,0.9,0.9}{
            \makecell[{{p{\linewidth}}}]{
                \texttt{\tiny{[GM$|$GM]}}
                \texttt{attempts = 1} \\
\texttt{target\_word = total} \\
\texttt{guess = arroz} \\
\texttt{guess\_feedback = a<yellow> r<red> r<red> o<yellow> z<red>} \\
            }
        }
    }
    & & \\ \\

    \theutterance \stepcounter{utterance}  
    & & & \multicolumn{4}{p{0.6\linewidth}}{
        \cellcolor[rgb]{0.9,0.9,0.9}{
            \makecell[{{p{\linewidth}}}]{
                \texttt{\tiny{[P1$\langle$GM]}}
                \texttt{retroalimentación\_intento: a<amarillo> r<rojo> r<rojo> o<amarillo> z<rojo>} \\
\\ 
\texttt{Proporciona tu respuesta solo en este formato.} \\
\texttt{explicación: detalles} \\
\texttt{intento: palabra} \\
            }
        }
    }
     \\ \\

    \theutterance \stepcounter{utterance}  
    & \multicolumn{4}{p{0.6\linewidth}}{
        \cellcolor[rgb]{0.561,0.737,0.561}{
            \makecell[{{p{\linewidth}}}]{
                \texttt{\tiny{[P1$\rangle$GM]}}
                \texttt{explicación: La retroalimentación indica que la palabra contiene una 'a' y una 'o' (amarillo), pero no 'r' ni 'z' (rojo). La 'a' no está en posición 1 y la 'o' no está en posición 4. Uso "palo" como base (con 'a' en pos 2 y 'o' en pos 4) y añado 'n' y 'e' (letras muy frecuentes) para probar "panel", colocando 'a' en pos 2 y 'o' en pos 5 (cumpliendo restricciones de posición).} \\
\texttt{intento: panel} \\
            }
        }
    }
    & & \\ \\

    \theutterance \stepcounter{utterance}  
    & & & \multicolumn{2}{p{0.3\linewidth}}{
        \cellcolor[rgb]{0.9,0.9,0.9}{
            \makecell[{{p{\linewidth}}}]{
                \texttt{\tiny{[GM$|$GM]}}
                \texttt{attempts = 2} \\
\texttt{target\_word = total} \\
\texttt{guess = panel} \\
\texttt{guess\_feedback = p<red> a<yellow> n<red> e<red> l<green>} \\
            }
        }
    }
    & & \\ \\

    \theutterance \stepcounter{utterance}  
    & & & \multicolumn{4}{p{0.6\linewidth}}{
        \cellcolor[rgb]{0.9,0.9,0.9}{
            \makecell[{{p{\linewidth}}}]{
                \texttt{\tiny{[P1$\langle$GM]}}
                \texttt{retroalimentación\_intento: p<rojo> a<amarillo> n<rojo> e<rojo> l<verde>} \\
\\ 
\texttt{Proporciona tu respuesta solo en este formato.} \\
\texttt{explicación: detalles} \\
\texttt{intento: palabra} \\
            }
        }
    }
     \\ \\

    \theutterance \stepcounter{utterance}  
    & \multicolumn{4}{p{0.6\linewidth}}{
        \cellcolor[rgb]{0.561,0.737,0.561}{
            \makecell[{{p{\linewidth}}}]{
                \texttt{\tiny{[P1$\rangle$GM]}}
                \texttt{explicación: La 'l' es verde en posición 5. La 'a' es amarilla en posición 2, por lo que no está en 1, 2 ni 5; debe ir en 3 o 4. La 'o' (amarilla en intento 1) no está en 4, y al no poder estar en 1, 2 ni 5, debe ir obligatoriamente en la 3. Por tanto, la 'a' va en la 4. Las letras p, n, e, r, z están descartadas. Pruebo consonantes frecuentes restantes (s, c, m, d) en posiciones 1 y 2. "saco" no sirve (a en 2). "sacol" cumple: s(1), a(4), c(2), o(3), l(5).} \\
\texttt{intento: sacol} \\
            }
        }
    }
    & & \\ \\

    \theutterance \stepcounter{utterance}  
    & & & \multicolumn{2}{p{0.3\linewidth}}{
        \cellcolor[rgb]{0.9,0.9,0.9}{
            \makecell[{{p{\linewidth}}}]{
                \texttt{\tiny{[GM$|$GM]}}
                \texttt{Error: The guessed word is not a valid word for this game.} \\
            }
        }
    }
    & & \\ \\

    \theutterance \stepcounter{utterance}  
    & & & \multicolumn{4}{p{0.6\linewidth}}{
        \cellcolor[rgb]{0.9,0.9,0.9}{
            \makecell[{{p{\linewidth}}}]{
                \texttt{\tiny{[P1$\langle$GM]}}
                \texttt{Tu intento no es una palabra válida para este juego. Por favor, inténtalo de nuevo.} \\
\\ 
\texttt{Proporciona tu respuesta solo en este formato.} \\
\texttt{explicación: detalles} \\
\texttt{intento: palabra} \\
            }
        }
    }
     \\ \\

    \theutterance \stepcounter{utterance}  
    & \multicolumn{4}{p{0.6\linewidth}}{
        \cellcolor[rgb]{0.561,0.737,0.561}{
            \makecell[{{p{\linewidth}}}]{
                \texttt{\tiny{[P1$\rangle$GM]}}
                \texttt{explicación: "sacol" no es válido. Mantengo el patrón deducido: \_ \_ o a l (o en 3, a en 4, l en 5). Pruebo "focal", palabra común que encaja en la estructura (f en 1, o en 3, c en 2, a en 4, l en 5) y usa letras nuevas frecuentes.} \\
\texttt{intento: focal} \\
            }
        }
    }
    & & \\ \\

    \theutterance \stepcounter{utterance}  
    & & & \multicolumn{2}{p{0.3\linewidth}}{
        \cellcolor[rgb]{0.9,0.9,0.9}{
            \makecell[{{p{\linewidth}}}]{
                \texttt{\tiny{[GM$|$GM]}}
                \texttt{attempts = 3} \\
\texttt{target\_word = total} \\
\texttt{guess = focal} \\
\texttt{guess\_feedback = f<red> o<green> c<red> a<green> l<green>} \\
            }
        }
    }
    & & \\ \\

    \theutterance \stepcounter{utterance}  
    & & & \multicolumn{4}{p{0.6\linewidth}}{
        \cellcolor[rgb]{0.9,0.9,0.9}{
            \makecell[{{p{\linewidth}}}]{
                \texttt{\tiny{[P1$\langle$GM]}}
                \texttt{retroalimentación\_intento: f<rojo> o<verde> c<rojo> a<verde> l<verde>} \\
\\ 
\texttt{Proporciona tu respuesta solo en este formato.} \\
\texttt{explicación: detalles} \\
\texttt{intento: palabra} \\
            }
        }
    }
     \\ \\

    \theutterance \stepcounter{utterance}  
    & \multicolumn{4}{p{0.6\linewidth}}{
        \cellcolor[rgb]{0.561,0.737,0.561}{
            \makecell[{{p{\linewidth}}}]{
                \texttt{\tiny{[P1$\rangle$GM]}}
                \texttt{explicación: Se confirma la estructura \_ \_ o a l con 'o', 'a', 'l' verdes en posiciones 3, 4, 5. Las letras f, c, p, n, e, r, z están descartadas (rojo). Pruebo 't' y 'r' (consonantes muy frecuentes) en posiciones 1 y 2 formando "troll", palabra válida que encaja en el patrón.} \\
\texttt{intento: troll} \\
            }
        }
    }
    & & \\ \\

    \theutterance \stepcounter{utterance}  
    & & & \multicolumn{2}{p{0.3\linewidth}}{
        \cellcolor[rgb]{0.9,0.9,0.9}{
            \makecell[{{p{\linewidth}}}]{
                \texttt{\tiny{[GM$|$GM]}}
                \texttt{Error: The guessed word is not a valid word for this game.} \\
            }
        }
    }
    & & \\ \\

    \theutterance \stepcounter{utterance}  
    & & & \multicolumn{4}{p{0.6\linewidth}}{
        \cellcolor[rgb]{0.9,0.9,0.9}{
            \makecell[{{p{\linewidth}}}]{
                \texttt{\tiny{[P1$\langle$GM]}}
                \texttt{Tu intento no es una palabra válida para este juego. Por favor, inténtalo de nuevo.} \\
\\ 
\texttt{Proporciona tu respuesta solo en este formato.} \\
\texttt{explicación: detalles} \\
\texttt{intento: palabra} \\
            }
        }
    }
     \\ \\

    \theutterance \stepcounter{utterance}  
    & \multicolumn{4}{p{0.6\linewidth}}{
        \cellcolor[rgb]{0.561,0.737,0.561}{
            \makecell[{{p{\linewidth}}}]{
                \texttt{\tiny{[P1$\rangle$GM]}}
                \texttt{explicación: "troll" no está en el diccionario del juego. Mantengo el patrón \_ \_ o a l. Pruebo "vocal", palabra muy común que encaja perfectamente (v en 1, o en 3, c en 2, a en 4, l en 5) y verifica la 'c' (aún no probada en esta posición) y la 'v'.} \\
\texttt{intento: vocal} \\
            }
        }
    }
    & & \\ \\

    \theutterance \stepcounter{utterance}  
    & & & \multicolumn{2}{p{0.3\linewidth}}{
        \cellcolor[rgb]{0.9,0.9,0.9}{
            \makecell[{{p{\linewidth}}}]{
                \texttt{\tiny{[GM$|$GM]}}
                \texttt{attempts = 4} \\
\texttt{target\_word = total} \\
\texttt{guess = vocal} \\
\texttt{guess\_feedback = v<red> o<green> c<red> a<green> l<green>} \\
            }
        }
    }
    & & \\ \\

    \theutterance \stepcounter{utterance}  
    & & & \multicolumn{4}{p{0.6\linewidth}}{
        \cellcolor[rgb]{0.9,0.9,0.9}{
            \makecell[{{p{\linewidth}}}]{
                \texttt{\tiny{[P1$\langle$GM]}}
                \texttt{retroalimentación\_intento: v<rojo> o<verde> c<rojo> a<verde> l<verde>} \\
\\ 
\texttt{Proporciona tu respuesta solo en este formato.} \\
\texttt{explicación: detalles} \\
\texttt{intento: palabra} \\
            }
        }
    }
     \\ \\

    \theutterance \stepcounter{utterance}  
    & \multicolumn{4}{p{0.6\linewidth}}{
        \cellcolor[rgb]{0.561,0.737,0.561}{
            \makecell[{{p{\linewidth}}}]{
                \texttt{\tiny{[P1$\rangle$GM]}}
                \texttt{explicación: Se confirma \_ \_ o a l. Letras descartadas: v, c, f, p, n, e, r, z. Quedan consonantes frecuentes para posiciones 1 y 2: s, t, d, m, b, g. "s" y "t" son las más probables. "s" en 1 y "t" en 2 forma "st" (inicio muy común). "stoal" no es palabra. "t" en 1 y "s" en 2 forma "ts" (poco común). Pruebo "total" (t, o, t, a, l), palabra muy frecuente que encaja y testa 't' en 1 y 3 (3 ya es 'o', así que 't' iría en 1).} \\
\texttt{intento: total} \\
            }
        }
    }
    & & \\ \\

    \theutterance \stepcounter{utterance}  
    & & & \multicolumn{2}{p{0.3\linewidth}}{
        \cellcolor[rgb]{0.9,0.9,0.9}{
            \makecell[{{p{\linewidth}}}]{
                \texttt{\tiny{[GM$|$GM]}}
                \texttt{attempts = 5} \\
\texttt{target\_word = total} \\
\texttt{guess = total} \\
\texttt{guess\_feedback = t<green> o<green> t<green> a<green> l<green>} \\
            }
        }
    }
    & & \\ \\

    \theutterance \stepcounter{utterance}  
    & & & \multicolumn{2}{p{0.3\linewidth}}{
        \cellcolor[rgb]{0.9,0.9,0.9}{
            \makecell[{{p{\linewidth}}}]{
                \texttt{\tiny{[GM$|$GM]}}
                \texttt{game\_result = WIN} \\
            }
        }
    }
    & & \\ \\

\end{supertabular}
}

\end{document}
